% ============================================================
% Chapter 6: Experimental Design
%   Section: Evaluation Harness
%     Subsections: The Tuning Loop; Run-to-Run Consistency
% Included from chapters/06-experimental-design/04-evaluation-harness/evaluation-harness.tex
% Sources: backend/src/evaluation/runner.py (HarnessRunner._loop),
% loop.py, consistency.py; docs/decisions/evaluation.md, "The second
% look", "Consistency".
% ============================================================

\subsection{The Tuning Loop}
\label{subsec:tuning-loop}

On the tuning half of the custom set, and only there, every claim goes
through a loop after its first verdict (\texttt{run --loop}). The loop
asks of each disagreement with the gold label whether a second attempt
fixes it, and if not, whether the claim or the system is at fault:

\begin{enumerate}
    \item \textbf{Verify.} The claim is checked as the pipeline checks
    it. If the verdict equals the gold label, and the pipeline's own
    trigger would not fire, the claim is done.
    \item \textbf{Look again.} Otherwise the claim gets a second look
    (Subsection~\ref{subsec:second-look}). If the second verdict equals
    the gold label, the claim is \emph{fixed}: the evidence was on the
    web, and the first search did not find it.
    \item \textbf{Check the claim itself.} If it still disagrees, the
    claim is judged once more, with no search, against the annotator's
    own evidence links --- the pages the gold label was set from, the
    ``oracle'' evidence condition of FEVER and AVeriTeC
    \cite{thorne2018fever,schlichtkrull2023averitec}. If on those pages
    the model reaches the gold label, the error is the
    \emph{retrieval}'s: the decisive source exists and the search never
    found it. If it gives the same verdict the pipeline gave, the model
    and the label disagree on the same evidence, and a person decides
    which is wrong (\emph{adjudicate}); a third verdict is a
    \emph{reasoning} error.
\end{enumerate}

The outcomes are reported with the claims behind each, and the
\emph{adjudicate} list goes to the annotator's review of the AI-assisted
labels (Subsection~\ref{subsec:custom-procedure}): a gold set an LLM
helped to label is the first thing the loop can find wrong. The loop
uses the gold label to choose which claims to look at again, so its
accuracy is an upper bound, not a score. What decides whether the second
look is turned on is the other number it reports: the accuracy obtained
when the second verdict replaces the first wherever the pipeline's own
trigger fires, without gold, paired over the same claims, with the
claims it fixes and the claims it breaks. Only that rule, chosen on the
tuning half, is run on the test half.

\subsection{Run-to-Run Consistency}
\label{subsec:consistency}

The fifth of the test half that the annotator labels twice
(Subsection~\ref{subsec:custom-review}) is also run through the system
twice, with the same model, thresholds and corpus snapshot, at a later
time. The two runs are compared over the claims both scored:

\begin{itemize}
    \item \textbf{Verdict agreement and Cohen's $\kappa$} between the
    runs, with percentile bootstrap intervals, and the confusion of one
    run's verdicts by the other's --- on the scale of the annotator's
    self-agreement, so the two can be read side by side; and the same
    for the model's raw verdict before recalibration.
    \item \textbf{Where a change came from}: the overlap of the two
    runs' ranked evidence for each claim (Jaccard over normalised
    addresses). A verdict that changed on the very same evidence is the
    model's own variation; one that changed with the evidence comes from
    the web, which changed between the two runs.
    \item \textbf{Whether it matters}: each run's accuracy against the
    gold labels on the same claims, and the paired difference --- two
    runs can disagree on a claim and both be wrong.
\end{itemize}
